\begin{lemma}[Infinite extension of a finite prefix]
Every finite set $s\subseteq\mathbb N$ is a proper initial segment of the
range of some increasing enumeration of an infinite subset of $\mathbb N$.
\end{lemma}
\begin{proof}
We argue separately according to whether $s$ is empty.

Suppose first that $s$ is nonempty.  Write its elements in increasing order
as
\[
a_0<a_1<\cdots<a_{r-1},
\]
where $r=|s|$, and let $m=a_{r-1}=\max s$.  Define $x:\mathbb N\to\mathbb
N$ by
\[
x(i)=a_i\quad (i<r),
\qquad
x(i)=m+1+(i-r)\quad (i\ge r).
\]
The first part is strictly increasing by the displayed ordering of the
elements of $s$.  The last element of the first part is $m$, while the first
element of the second part is $m+1$, so the inequality also holds across the
boundary.  For $r\le i<j$, the formula for the second part gives
$x(i)<x(j)$.  Thus $x$ is an order embedding, hence an element of
$\mathsf{IncSeq}$ by \noderef{IncSeq}.

The first $r$ values of $x$ are exactly the elements of $s$, and the
remaining values are exactly $m+1,m+2,\ldots$.  Consequently
\[
\operatorname{range}(x)=s\cup\{m+1,m+2,\ldots\}.
\]
In particular, $m+1\in\operatorname{range}(x)$, and no value from the tail is
less than $m+1$; every element of $s$ is at most $m$ and therefore is less
than $m+1$.  Hence
\[
s=\{k\in\operatorname{range}(x)\mid k<m+1\}.
\]
This exhibits $m+1$ as the required proper-prefix witness in the definition
of $\noderef{ProperPrefixSet}$.

Now suppose that $s=\varnothing$.  Use the identity map $x(i)=i$.  If
$i<j$, then $x(i)=i<j=x(j)$, so $x$ is an order embedding and hence belongs
to $\mathsf{IncSeq}$ by \noderef{IncSeq}.  Its range is all of $\mathbb N$.
The number $0$ belongs to this range, and
\[
\varnothing=\{k\in\mathbb N\mid k<0\}
       =\{k\in\operatorname{range}(x)\mid k<0\}.
\]
Thus $0$ is a proper-prefix witness, by \noderef{ProperPrefixSet}.
\end{proof}
